\documentclass[9pt]{extarticle}

% ===== Page geometry: one-page poster, tight margins =====
\usepackage[a4paper, landscape, top=0.3in, bottom=0.2in, left=0.35in, right=0.35in]{geometry}

% ===== Core packages =====
\usepackage[utf8]{inputenc}
\usepackage[T1]{fontenc}
\usepackage{multicol}
\usepackage{booktabs}
\usepackage{array}
\usepackage{tabularx}
\usepackage{colortbl}
\usepackage{multirow}
\usepackage{xcolor}
\usepackage[final]{graphicx}
\usepackage{titlesec}
\usepackage[most]{tcolorbox}
\usepackage{enumitem}
\usepackage[expansion=false]{microtype}
\usepackage{wrapfig}

% ===== Colors =====
\definecolor{primary}{RGB}{25,70,130}
\definecolor{accent}{RGB}{200,70,40}
\definecolor{lightgray}{RGB}{235,238,245}
\definecolor{softyellow}{RGB}{255,248,220}
\definecolor{darkgray}{RGB}{70,70,70}
\definecolor{successgreen}{RGB}{40,130,70}
\definecolor{heroblue}{RGB}{220,234,255}

% ===== Section style =====
\titleformat{\section}
  {\normalfont\normalsize\bfseries\color{primary}}
  {\thesection.}{0.3em}{}[\vspace{-0.3em}{\color{primary}\hrule height 0.5pt}]
\titlespacing*{\section}{0pt}{0.2em}{0.05em}

% ===== No page numbers =====
\pagestyle{empty}
\setlength{\columnsep}{0.8em}
\setlength{\parindent}{0pt}
\setlength{\parskip}{0.1em}

% ===== Custom boxes =====
\newtcolorbox{takeawaybox}{
  colback=lightgray, colframe=primary, boxrule=0.6pt,
  arc=2pt, left=5pt, right=5pt, top=2pt, bottom=2pt
}

\newtcolorbox{heroboxnote}{
  colback=softyellow, colframe=accent, boxrule=0.6pt,
  arc=2pt, left=5pt, right=5pt, top=2pt, bottom=2pt
}

\newtcolorbox{herobluebox}{
  colback=heroblue, colframe=primary, boxrule=0.8pt,
  arc=2pt, left=6pt, right=6pt, top=2pt, bottom=2pt
}

\newtcolorbox{finding}{
  colback=white, colframe=primary, boxrule=0.4pt,
  arc=2pt, left=4pt, right=4pt, top=2pt, bottom=2pt
}

% ===== Compact list spacing =====
\setlist[itemize]{leftmargin=1.1em, itemsep=0.03em, topsep=0.03em, parsep=0pt}
\setlist[enumerate]{leftmargin=1.3em, itemsep=0.05em, topsep=0.03em, parsep=0pt}

% ===== Table spacing =====
\setlength{\aboverulesep}{0pt}
\setlength{\belowrulesep}{0pt}
\setlength{\extrarowheight}{0.3pt}
\setlength{\tabcolsep}{3pt}

\begin{document}

% ============ TITLE BLOCK ============
\begin{center}
  {\LARGE\bfseries\color{primary} Joint Alignment Speculative Decoding for Real-Time LLM Role-Play}\\[0.05em]
  {\footnotesize\itshape\color{darkgray} A draft model aligned with the target's \textbf{deployment distribution}, not its fine-tuning distribution}\quad
  {\footnotesize Siqi Zhang $\bullet$ CSE-398 $\bullet$ Spring 2026}
\end{center}
\vspace{-0.6em}
{\color{primary}\hrule height 1pt}
\vspace{0.1em}

% ============ HEADLINE RESULT BANNER ============
\begin{herobluebox}
\small\textbf{Headline.} Our \emph{Joint Alignment} draft training is the \textbf{only} configuration with higher acceptance than the off-the-shelf baseline on \textbf{both} in-character (+4.4\%) \textbf{and} off-topic (+6.3\%) queries. The naive ``train-on-character-data-only'' alternative (B3) is fastest on Jack Sparrow prompts but \textbf{silently loses 15.1\%} acceleration on off-topic prompts --- a failure mode invisible to in-domain-only evaluation.
\end{herobluebox}

% ============ MAIN 3-COLUMN LAYOUT ============
\begin{multicols}{3}

% ---------- COLUMN 1 ----------
\section{Problem \& Motivation}

Deploying a persona-specialized LLM (Jack Sparrow chatbot) requires \textbf{real-time speed} ($\geq$40 tok/s) \emph{and} \textbf{persona fidelity}. LoRA fine-tuning the target gives fidelity; \textbf{speculative decoding} (EAGLE 3) gives speed via a small \emph{draft} model that proposes tokens the target verifies in parallel.

Because rejection sampling preserves the target's exact output distribution, \textbf{the draft is a pure speed dial --- it cannot degrade quality}. So the question becomes: \emph{which draft maximizes acceleration in realistic deployment?}

\textbf{Open question.} Given a persona-specialized target, should the draft be persona-specialized too? Persona-only training looks fastest on in-character queries, but a deployed chatbot still receives off-topic user input. Does narrow draft training cost speed there?

\section{Method: Joint Alignment}

\begin{heroboxnote}
\textbf{Core insight.} Real users may send both in-character and off-topic queries.
\end{heroboxnote}

\textbf{Our recipe.} Training the EAGLE 3 draft model against the persona-tuned target's logits via a 7-step KL-divergence loss, on a balanced corpus:
\begin{itemize}
  \item \textbf{50\% persona data}: 1{,}429 Jack Sparrow samples (RoleBench).
  \item \textbf{50\% general data}: 1{,}429 UltraChat 200K samples.
\end{itemize}
Train for \textbf{8 epochs}.

\textbf{The configurations under comparison.} All three share the same persona-tuned target: \textbf{Qwen3-4B} + LoRA on RoleBench Jack Sparrow; they differ only in the + draft model (EAGLE 3 head for Qwen3-4B, $\sim$0.2B params): :

\small
\begin{tabularx}{\linewidth}{@{}lX@{}}
\toprule
\textbf{B2$'$} & off-the-shelf EAGLE 3 draft \emph{(no retraining --- the draft was never aligned to this target)} \\
\textbf{B3}    & draft retrained on \textbf{100\% Jack Sparrow} data, 20 epochs \emph{(naive specialization)} \\
\rowcolor{lightgray}\textbf{Prop-B (ours)} & draft retrained on the \textbf{50/50 mix} above, 8 epochs \emph{(Joint Alignment)} \\
\bottomrule
\end{tabularx}
\normalsize

\section{Novelty}

Prior persona work has rarely adopted speculative decoding. We bring a state-of-the-art acceleration algorithm into the role-play setting, reduce the amount of waiting time for users. Crucially, our \emph{Joint Alignment} draft preserves acceleration on \textbf{off-topic} user queries as well as in-character ones, giving the full pipeline stronger generalization.

\section{Experimental Setup}

\textbf{Two disjoint test sets} (no training-data overlap):
\begin{itemize}
  \item \textbf{In-domain:} RoleBench Jack Sparrow test (349 prompts + references).
  \item \textbf{Out-of-domain:} UltraChat 200K (349--1{,}000 prompts).
\end{itemize}

\textbf{Three metric families:} \textit{Acceptance length} (how many tokens are accepted by target model), \textit{TPS} (fixed 256-token output), \textit{BLEU/ROUGE-L} (quality check; verifies lossless guarantee).

% ---------- COLUMN 2 ----------

\section{Speedup Comparison \normalfont\small (TPS, 5 configs)}

\small
\begin{tabularx}{\linewidth}{@{}lXcccc@{}}
\toprule
\textbf{Cfg} & \textbf{Target / Draft} & \textbf{JS} & \textbf{UC} & \textbf{JS$\times$} & \textbf{UC$\times$} \\
\midrule
B1     & Base / ---               & 24.45 & 23.13 & 1.00 & 1.00 \\
B2     & Base / general           & 45.21 & ---   & 1.85 & ---  \\
B2$'$  & JS-FT / general          & 62.03 & 55.39 & 2.54 & 2.39 \\
B3     & JS-FT / JS-only          & \textbf{66.89} & 55.09 & \textbf{2.74} & 2.38 \\
\rowcolor{lightgray}\textbf{Prop-B} & \textbf{JS-FT / mix (ours)} & 58.75 & \textbf{58.86} & 2.40 & \textbf{2.54} \\
\bottomrule
\end{tabularx}
\normalsize
\vspace{0.05em}\\
\footnotesize\emph{Speedup vs domain-matched B1 baseline (JS = 24.45 TPS, UC = 23.13 TPS). B3 is fastest on JS but its OOD speedup (2.38$\times$) is the \textbf{lowest} of the speculative configs --- its specialization gain does not transfer. Prop-B is the only config with UC speedup $>$ 2.50.}
\normalsize

\section{Acceptance Length \normalfont\small (draft--target alignment)}

\small
\begin{tabularx}{\linewidth}{@{}lXccX@{}}
\toprule
\textbf{Cfg} & \textbf{Draft} & \textbf{JS} & \textbf{UC} & \textbf{UC $\Delta$} \\
\midrule
B2$'$  & off-the-shelf    & 3.84 & 3.63 & \textcolor{darkgray}{$-5.5\%$} \\
B3     & JS-only          & \textbf{4.25} & 3.61 & \textcolor{accent}{\textbf{$-15.1\%$}} \\
\rowcolor{lightgray}\textbf{Prop-B} & \textbf{mix (ours)} & \textbf{4.01} & \textbf{3.86} & \textcolor{successgreen}{\textbf{$-3.7\%$}} \\
\bottomrule
\end{tabularx}
\normalsize
\vspace{0.05em}\\
\footnotesize\emph{Prop-B is the only config higher than B2$'$ on \textbf{both} domains (+4.4\% JS, +6.3\% UC). B3 pays 15.1\% on UC for its JS peak.}
\normalsize

\vspace{-0.15em}
\begin{center}
\includegraphics[width=0.95\linewidth,height=4.2cm,keepaspectratio]{pareto_plot.png}
\end{center}
\vspace{-0.95em}
\footnotesize Upper-right = better. Prop-B (star) is the \emph{only} point that dominates the B2$'$ baseline on both axes.
\normalsize

\vspace{-0.4em}
\section{Quality (Reference Similarity)}

BLEU + ROUGE-L vs ground-truth, free-length generation:

\small
\begin{tabularx}{\linewidth}{@{}XlXcXc@{}}
\toprule
\multirow{2}{*}{\textbf{Cfg}} & \multirow{2}{*}{\textbf{Target/Draft}} & \multicolumn{2}{c}{\textbf{JS}} & \multicolumn{2}{c}{\textbf{UC}} \\
\cmidrule(lr){3-4}\cmidrule(lr){5-6}
       &                 & BLEU & R-L   & BLEU & R-L \\
\midrule
B1     & Base / AR       & 2.84 & 15.52 & 7.23 & 24.37 \\
B3     & JS-FT / JS      & \textbf{9.06} & \textbf{24.43} & 10.98 & 24.39 \\
\rowcolor{lightgray}Prop-B & JS-FT / mix & 8.52 & 23.99 & \textbf{11.30} & \textbf{25.16} \\
\bottomrule
\end{tabularx}
\normalsize

\vspace{0.1em}
\footnotesize\emph{Both finetuned configs crush B1 (JS BLEU $\sim3\times$). Prop-B matches B3 in-domain (spread $<$2\%, sampling noise), exceeds B3 on OOD. Speculative decoding is \textbf{mathematically lossless} $\Rightarrow$ all JS-FT configs share the same output distribution; draft training is a \textbf{pure speed-axis} optimization.}
\normalsize

% ---------- COLUMN 3 ----------

\section{Why Prop-B Outperforms}

\begin{takeawaybox}
\begin{enumerate}
  \item \textbf{Balanced speed:} only config $>$55 TPS on \emph{both} JS and UC (58.75 / 58.86). B3 drops to 55.09 on UC.
  \item \textbf{Pareto-dominates B2$'$ on acceptance:} +4.4\% in-domain \emph{and} +6.3\% out-of-domain. B3 dominates only in-domain.
  \item \textbf{Highest OOD acceptance overall:} 3.86 vs B3's 3.61 (+\textbf{6.9\%}) and B2$'$'s 3.63 (+\textbf{6.3\%}).
  \item \textbf{Lossless quality:} matches or exceeds B3 on BLEU / ROUGE-L across both domains.
  \item \textbf{Training-data efficient:} single-stage 8-epoch training on 2.9k-sample mixed corpus.
\end{enumerate}
\end{takeawaybox}

\footnotesize\emph{Ablation supporting recipe choices: (i) 22\% $\to$ 50\% character share improves JS acceptance $3.84 \to 4.01$; (ii) training 20 epochs instead of 8 regresses the draft on both axes (3.53 / 3.63) despite lower train loss --- overfitting to the training mix.}
\normalsize

\section{Conclusions}

\begin{enumerate}
  \item \textbf{Naive persona-only draft training silently degrades $\sim$15\% of OOD deployment speed}.
  \item \textbf{Joint Alignment} (50/50 persona + general, 8 epochs) is the only configuration that Pareto-dominates the off-shelf baseline on both axes, delivering $\mathbf{2.4\times}$ speedup with \emph{no} quality cost.
  \item \textbf{Under speculative decoding, speed and quality are decoupled}: train the draft for whatever distribution reflects deployment; the target fixes output quality.
\end{enumerate}

\textbf{Limitations.} Single persona (Jack Sparrow); limited training data; EAGLE 3 only, where other algorithms such as DFlash may behave differently.

\end{multicols}

\vspace{-0.7em}
\begin{center}
\scriptsize
\textcolor{primary}{\rule{0.9\linewidth}{0.4pt}}\\[0.05em]
\textbf{Code \& full hyperparameters:} \texttt{github.com/...} \quad$\bullet$\quad
\textbf{Hardware:} 4$\times$ RTX 2080 Ti
\end{center}

\end{document}